% Source: https://artofproblemsolving.com/community/c3h1568552_mc_handbook
% Stable identity: aops-topic-1568552-q02
Each dimension occurs on $4$ edges, so the edges add up to $4(a+b+c)=12\cdot13$, and
\[
a+b+c=39,\qquad a<b<c,\qquad b^2=ac .
\]
\textbf{Parametrising.} Write the common ratio $\dfrac {b}{a}=\dfrac {c}{b}$ in lowest terms as $\dfrac {q}{p}$ with $q>p\ge1$. From $b=\dfrac{aq}{p}$ we get $p\mid a$, and from $c=\dfrac{aq^2}{p^2}$, $p^2\mid aq^2$, so $p^2\mid a$ (as $\gcd(p,q)=1$). With $a=kp^2$:
\[
a=kp^2,\qquad b=kpq,\qquad c=kq^2,\qquad k\left(p^2+pq+q^2\right)=39 .
\]
\textbf{Cases.} The value $p^2+pq+q^2$ is at least $1+2+4=7$ and divides $39=3\cdot13$, so it is $13$ or $39$.
\begin{itemize}
\item $p^2+pq+q^2=13$: only $(p,q)=(1,3)$, with $k=3$: dimensions $3,9,27$.
\item $p^2+pq+q^2=39$: only $(p,q)=(2,5)$, with $k=1$: dimensions $4,10,25$. (Checking $p=1,2,3$: $1+q+q^2=39$ and $9+3q+q^2=39$ have no integer roots.)
\end{itemize}

\begin{asy}
usepackage("amsmath");
size(6cm);
defaultpen(fontsize(10pt));
void box(pair o,real a,real b,real c,pen p){
  pair k=(0.42,0.3)*1.0;
  pair p0=o,p1=o+(a,0),p2=o+(a,0)+b*k,p3=o+b*k;pair u=(0,c);
  draw(p0--p1--p1+u--p0+u--cycle,p);draw(p1--p2--p2+u--p1+u,p);draw(p0+u--p3+u--p2+u,p);
  draw(p0--p3,dashed);draw(p3--p2,dashed);draw(p3--p3+u,dashed);
}
real s=0.25;
box((0,0),27*s,9*s,3*s,linewidth(0.9));
label("$27$",(27*s/2,0),S);label("$9$",(27*s,0)+(9*s/2)*(0.42,0.3),SE);label("$3$",(0,1.5*s),W);
box((0,-4.5),25*s,10*s,4*s,linewidth(0.9));
label("$25$",(25*s/2,-4.5),S);label("$10$",(25*s,-4.5)+(10*s/2)*(0.42,0.3),SE);label("$4$",(0,-4.5+2*s),W);
label("\mbox{}",(min(currentpicture,user=true)+max(currentpicture,user=true))/2,black); // keeps the text after the figure black when LaTeX typesets the labels inline
\end{asy}

\textbf{Volumes.} For dimensions in geometric progression the volume is $abc=b^3$. So the volumes are $9^3=729$ and $10^3=1000$, and their sum is
\[
729+1000=\boxed{1729}.
\]
